\documentclass[11pt,a4paper]{article}

\usepackage[margin=2.2cm]{geometry}
\usepackage[T1]{fontenc}
\usepackage{lmodern}
\usepackage{microtype}
\usepackage{amsmath}
\usepackage{booktabs}
\usepackage{longtable}
\usepackage{array}
\usepackage{xcolor}
\usepackage{listings}
\usepackage[hidelinks]{hyperref}

\definecolor{codebg}{RGB}{246,248,250}
\definecolor{coderule}{RGB}{208,215,222}
\definecolor{accent}{RGB}{31,93,140}

\lstset{
  basicstyle=\ttfamily\small,
  backgroundcolor=\color{codebg},
  frame=single,
  rulecolor=\color{coderule},
  breaklines=true,
  columns=fullflexible,
  keepspaces=true,
  showstringspaces=false,
  xleftmargin=0.3em,
  xrightmargin=0.3em
}

\newcommand{\code}[1]{\texttt{#1}}
\newcommand{\sourcepath}[1]{\par\smallskip\noindent\textit{Source:} \path{#1}\par\smallskip}

\title{\textbf{Lab 1 Report: Running Synthetic Traffic}\\[0.4em]
       \large gem5 Garnet Network-on-Chip Simulation}
\author{\begin{tabular}{rl}
          Name: & \rule{7cm}{0.4pt} \\
          Student ID: & \rule{7cm}{0.4pt}
        \end{tabular}}
\date{gem5 v23.0.0.1}

\begin{document}
\maketitle

\begin{abstract}
This laboratory studies the Garnet 3.0 network model in gem5 using synthetic
traffic. A 64-node, $8\times8$ mesh is exercised with uniform-random,
single-flit control traffic. The global tick frequency is corrected to
2\,GHz, statistical units are added to Garnet, and the measured packet rate,
latency, and hop count are analyzed. The final run injects 3165 packets and
receives 3162 packets, with a reception rate of
$0.004940625$ packets/node/network-cycle.
\end{abstract}

\tableofcontents
\newpage

\section{Experimental Setup}

The experiment uses gem5 tag \code{v23.0.0.1}, commit
\code{af72b9ba580546ac12ce05bfaac3fd53fa8699f4}, built for the NULL ISA with
the \code{Garnet\_standalone} protocol. The relevant incremental build command
is:

\begin{lstlisting}[language=bash]
scons --ignore-style build/NULL/gem5.opt \
  PROTOCOL=Garnet_standalone -j$(nproc)
\end{lstlisting}

SCons tracks dependencies, so changing \code{GarnetNetwork.cc} recompiles the
affected Garnet object and relinks \code{gem5.opt}; it does not require deleting
the build directory or recompiling every source file.

The test configuration is:

\begin{lstlisting}[language=bash]
./build/NULL/gem5.opt \
  configs/example/garnet_synth_traffic.py \
  --network=garnet --num-cpus=64 --num-dirs=64 \
  --topology=Mesh_XY --mesh-rows=8 \
  --inj-vnet=0 --synthetic=uniform_random \
  --sim-cycles=10000 --injectionrate=0.01
\end{lstlisting}

The 64 objects named CPUs are synthetic traffic generators rather than
instruction-executing processors. \code{Mesh\_XY} creates 64 routers arranged
as eight rows by eight columns. There are also 64 directory destinations.
Virtual network 0 carries one-flit control packets, and every source chooses a
destination uniformly from all 64 destinations, including itself.

\section{Task 1: Frequency Correction and Test Run}

\subsection{Global frequency change}

The original configuration used:

\begin{lstlisting}[language=Python]
m5.ticks.setGlobalFrequency("1ps")
\end{lstlisting}

This means one tick is 1\,ps, or $10^{12}$ ticks/s. It was changed to:

\begin{lstlisting}[language=Python]
m5.ticks.setGlobalFrequency("2GHz")
\end{lstlisting}

\sourcepath{/root/gem5/configs/example/garnet_synth_traffic.py}

The final statistics confirm
\code{simFreq = 2000000000} ticks/s and \code{simTicks = 10000}. Therefore,

\begin{equation}
T_{\mathrm{tick}}=\frac{1}{2\times10^9}=0.5\ \mathrm{ns},
\qquad
T_{\mathrm{sim}}=10000\times0.5\ \mathrm{ns}=5\ \mu\mathrm{s}.
\end{equation}

The system clock is 1\,GHz, while Ruby/Garnet runs at 2\,GHz. Consequently,

\begin{equation}
1\ \text{system cycle}=2\ \text{ticks}=1\ \mathrm{ns},\qquad
1\ \text{Ruby cycle}=1\ \text{tick}=0.5\ \mathrm{ns}.
\end{equation}

The run thus contains approximately 5000 traffic-generator cycles and 10000
Ruby network cycles.

\subsection{Measured results}

\begin{table}[h]
\centering
\caption{Statistics before and after correcting the global frequency.}
\begin{tabular}{lrr}
\toprule
Metric & Original 1\,ps/tick & Final 2\,GHz \\
\midrule
Packets injected & 5 & 3165 \\
Packets received & 2 & 3162 \\
Average queueing latency & 1000 ticks & 2 ticks/packet \\
Average network latency & 3000 ticks & 13.557559 ticks/packet \\
Average total latency & 4000 ticks & 15.557559 ticks/packet \\
Average hops & 1.500000 & 5.269450 \\
\bottomrule
\end{tabular}
\end{table}

At an injection probability of 0.01 per node per system cycle, the expected
number of generated packets is approximately

\begin{equation}
64\ \text{nodes}\times5000\ \text{cycles}\times0.01
=3200\ \text{packets}.
\end{equation}

The observed 3165 injections are consistent with this Bernoulli process. Only
three packets remain in flight at termination. For a uniform-random
$8\times8$ mesh, the expected Manhattan distance is

\begin{equation}
E[H]=2E[|X_s-X_d|]
=2\frac{8^2-1}{3\times8}=5.25\ \text{hops}.
\end{equation}

The measured 5.269450 hops differs from the theoretical value by only 0.37\%,
which is reasonable for 3162 received packets.

\section{Task 2: Statistical Units and Reception Rate}

The original Garnet statistics were emitted as \code{(Unspecified)}. Unit
metadata was added in \code{GarnetNetwork::regStats()} without changing any
counter or network behavior.

\begin{table}[h]
\centering
\caption{Units assigned to Garnet statistical variables.}
\begin{tabular}{>{\raggedright\arraybackslash}p{7.2cm}l}
\toprule
Statistics & Unit \\
\midrule
Injected/received packets and flits; traffic distribution & Count \\
Accumulated packet and flit latency & Tick \\
Average packet and flit latency & Tick/Count \\
Total hops and link activity & Count \\
Average hops & Count/Count \\
Average link utilization and VC load & Ratio \\
\bottomrule
\end{tabular}
\end{table}

\sourcepath{/root/gem5/src/mem/ruby/network/garnet/GarnetNetwork.cc}

The required reception rate is normalized by node count and the 10000 network
cycles specified by the laboratory:

\begin{equation}
R_{\mathrm{recv}}
=\frac{N_{\mathrm{received}}}{N_{\mathrm{nodes}}N_{\mathrm{cycles}}}
=\frac{3162}{64\times10000}
=0.004940625\ \frac{\mathrm{packets}}{\mathrm{node\cdot cycle}}.
\end{equation}

This is a throughput, not a delivery ratio. The delivery ratio is instead
$3162/3165\approx99.905\%$. The offered rate is
0.01 packets/node/system-cycle. Since one system cycle spans two Ruby cycles,
the equivalent offered network rate is 0.005 packets/node/Ruby-cycle, close to
the measured reception rate.

\section{Task 3: Configuration and Source-Code Analysis}

\subsection{Input parameters and defaults}

All accepted options can be displayed with:

\begin{lstlisting}[language=bash]
./build/NULL/gem5.opt configs/example/garnet_synth_traffic.py -h
\end{lstlisting}

The test-specific options are summarized below.

\begin{longtable}{@{}>{\ttfamily\small}p{3.5cm}p{3.0cm}p{7.6cm}@{}}
\caption{Synthetic traffic options.}\\
\toprule
\normalfont Option & \normalfont Default & \normalfont Purpose \\
\midrule
\endfirsthead
\toprule
\normalfont Option & \normalfont Default & \normalfont Purpose \\
\midrule
\endhead
--synthetic & \normalfont uniform\_random & \normalfont Destination pattern: uniform random, tornado, bit complement, bit reverse, bit rotation, neighbor, shuffle, or transpose. \\
--injectionrate & \normalfont 0.1 & \normalfont Packet-generation probability per node per traffic-generator cycle. \\
--precision & \normalfont 3 & \normalfont Decimal precision used by the random injection test. \\
--sim-cycles & \normalfont 1000 & \normalfont Simulation termination threshold; in this version it is compared directly with \code{curTick()}. \\
--num-packets-max & \normalfont -1 & \normalfont Maximum packets generated by each sender; -1 disables the limit. \\
--single-sender-id & \normalfont -1 & \normalfont Restrict injection to one sender; -1 enables every sender. \\
--single-dest-id & \normalfont -1 & \normalfont Force one destination; -1 uses the selected traffic pattern. \\
--inj-vnet & \normalfont -1 & \normalfont Select vnet 0, 1, or 2; -1 selects among them randomly. Vnets 0/1 use one-flit control packets and vnet 2 uses five-flit data packets. \\
\bottomrule
\end{longtable}

These defaults are defined directly in
\path{configs/example/garnet_synth_traffic.py}. General system and memory
options such as \code{--num-cpus}, \code{--num-dirs}, \code{--sys-clock},
memory type, cache geometry, and maximum tick come from
\path{configs/common/Options.py}. Ruby defaults, including the 2\,GHz Ruby
clock, come from \path{configs/ruby/Ruby.py}. Network defaults come from
\path{configs/network/Network.py}: Crossbar topology, simple network,
one-cycle router and link latency, 128-bit links, four VCs per vnet,
weight-table routing, and a 50000-cycle deadlock threshold.

For this experiment, \code{Mesh\_XY.py} assigns horizontal link weight 1 and
vertical link weight 2. The default weight-based routing table therefore
enforces dimension-ordered XY routing. Alternatively,
\code{--routing-algorithm=1} selects the explicit XY implementation in
\code{RoutingUnit.cc}.

\subsection{Units of cycles, router latency, and link latency}

A tick is the global, integer simulation time quantum. A cycle belongs to a
specific clock domain. Their relationship is

\begin{equation}
\text{ticks per component cycle}
=\frac{f_{\mathrm{global}}}{f_{\mathrm{component}}}.
\end{equation}

\code{router-latency} is a number of Ruby/router cycles and represents router
pipeline stages. An arriving flit is held in its input VC for
\code{router-latency - 1} cycles before switch allocation. Likewise,
\code{link-latency} is a number of network-link cycles; the link marks the flit
ready at \code{clockEdge(link\_latency)}.

Although the command-line help describes \code{sim-cycles} as cycles,
\path{GarnetSyntheticTraffic::tick()} compares \code{curTick()} directly with
\code{simCycles}. Its strict implementation unit in v23.0.0.1 is therefore
Tick. With both global and Ruby frequencies at 2\,GHz, one tick happens to equal
one Ruby cycle.

\subsection{Injection-rate unit}

Each \code{GarnetSyntheticTraffic} object wakes once per system clock cycle. It
draws a random integer and generates a request when the draw is below
\code{injRate} times the precision range. Thus the injection-rate unit is
packets/node/system-cycle. It is an offered probabilistic load, not a guarantee
that exactly that many packets enter or leave the network.

\subsection{Network interface and router definitions}

The Python SimObject declarations and their C++ implementations are:

\begin{table}[h]
\centering
\begin{tabular}{lll}
\toprule
SimObject & Python declaration & C++ implementation \\
\midrule
GarnetNetworkInterface & GarnetNetwork.py & NetworkInterface.hh/.cc \\
GarnetRouter & GarnetNetwork.py & Router.hh/.cc \\
\bottomrule
\end{tabular}
\end{table}

\noindent The files reside under
\path{src/mem/ruby/network/garnet/}. The configuration layer in
\path{configs/network/Network.py} selects these classes when
\code{--network=garnet} is used.

\subsection{Packet generation, buffering, and downstream eligibility}

\begin{center}
\fbox{\begin{minipage}{0.92\textwidth}\centering
GarnetSyntheticTraffic::generatePkt()
$\rightarrow$ RubyPort/Sequencer
$\rightarrow$ Garnet\_standalone cache MessageBuffer
$\rightarrow$ NetworkInterface::flitisizeMessage()
$\rightarrow$ source NI VCs
$\rightarrow$ NetworkLink
$\rightarrow$ Router InputUnit VC
$\rightarrow$ SwitchAllocator/Crossbar
$\rightarrow$ destination NI
$\rightarrow$ directory MessageBuffer
\end{minipage}}
\end{center}

The \code{generatePkt()} method in \code{GarnetSyntheticTraffic} chooses a destination and creates
a Read, instruction-fetch, or Write request. RubyPort and Sequencer deliver it
to the Garnet standalone cache controller. The controller maps requests to
vnet 0, 1, or 2 and enqueues a protocol message. The source NetworkInterface
allocates an output VC, converts the protocol message into flits, and increments
the injected statistics.

During transmission, flits may be stored in the protocol MessageBuffer, the
source NI's \code{niOutVcs}, a \code{NetworkLink} buffer, a router
\code{InputUnit}'s \code{VirtualChannel::inputBuffer}, a transient crossbar
buffer, or the destination NI stall queue. Packets are represented as flits
inside Garnet; the tail flit identifies packet completion.

\code{RoutingUnit} chooses the output direction but does not decide whether
the flit may advance. That decision is made primarily by
\code{SwitchAllocator::send\_allowed()}. A head flit needs a free downstream
output VC; a body or tail flit needs a positive credit, meaning an available
slot in the downstream input VC. Ordered vnets must also preserve the arrival
order of flits requesting the same output. \code{OutputUnit} and
\code{OutVcState} maintain credit counts, while reverse \code{CreditLink}s
return freed slots to the upstream router.

At the destination, the NetworkInterface places the completed protocol message
in the directory input buffer and records receive and latency statistics. The
Garnet standalone directory then dequeues and discards it because this protocol
is designed to measure network behavior rather than execute a complete cache
coherence transaction.

\section{Conclusion}

The corrected simulation runs for 10000 Ruby cycles and produces packet counts
consistent with the configured probabilistic load. The measured average hop
count agrees with the theoretical $8\times8$ mesh distance, and 99.905\% of
injected packets arrive before termination. Adding explicit units removes
ambiguity from Garnet's statistics, while source inspection shows that
credit-based flow control and switch allocation govern downstream movement.

\end{document}
